\chapter{Assembly Theory: Causal Depth and Copy Number}
\label{sec:assembly}
% ===================================================================

\section{Assembly Index as Causal Depth}

Assembly Theory assigns to each physical object $M$ an
\emph{assembly index} $\mathrm{AI}(M)$: the minimum number of
joining steps required to construct $M$ from copies of its
constituent parts, where each step joins two objects already
present into a new one.
The theory also proposes that \emph{time is an object}: the
assembly index records not a static property of $M$ but the
accumulated computational history of its construction.

This proposal maps precisely onto the notion of causal depth
developed in Part~I.

\begin{definition}[Elementary Terms]
Let $\GSLT$ be a GSLT.
A set $E \subseteq \terms(\GSLT)$ of \emph{elementary terms}
is a designated collection of starting materials --- the analogue
of atoms or monomers in chemistry --- from which all other terms
are constructed by rewriting.
\end{definition}

\begin{proposition}[Assembly Index as Minimum Causal Depth]
\label{prop:assembly-index}
The assembly index of a term $M \in \terms(\GSLT)$ with respect
to elementary terms $E$ is
\[
  \mathrm{AI}(M) \;=\; d^o_{\min}(E, M)
  \;=\; \min\bigl\{ |\gamma| \;\mid\;
    \gamma \text{ is a context-labeled path in } \mathcal{ST}_o
    \text{ from some } e \in E \text{ to } M
  \bigr\}.
\]
That is, $\mathrm{AI}(M)$ is the minimum path length in the
\emph{open} synchronization tree from the elementary terms to $M$,
where each edge is labeled by the minimal context required to
trigger the corresponding transition.
\end{proposition}

\begin{remark}[Why the Open Tree, Not the Closed Tree]
Assembly steps require an environment to supply a reaction partner ---
a catalyst, a monomer source, an energy input.
These are precisely the transitions that appear as edges in the
open synchronization tree (labeled by minimal contexts) rather than
the closed tree (which records only autonomous rewrites).
Each edge in $\mathcal{ST}_o$ labeled by a minimal context $K$
represents one joining step: the context $K$ is the environmental
participant that enables the construction.
The \emph{nesting depth} of context labels along a path is
therefore exactly the assembly index: each layer of context is
one more joining operation, one more step of accumulated causal
history.
This makes precise Cronin and Walker's intuition that assembly
index measures ``time as accumulating computational effects'':
the context labels \emph{are} the accumulated effects, and their
nesting depth \emph{is} the causal time elapsed.
The notion of time as an object in Assembly Theory is not a
metaphor; it is the minimum context-labeled path length in an
open synchronization tree.
\end{remark}

\section{The Two Gaps in Assembly Theory}
Despite the strength of the assembly index concept, Assembly Theory
has attracted criticism on two grounds.

\paragraph{Gap 1: Copy number lacks a notion of location.}
The \emph{copy number} $\mathrm{CN}(M)$ is supposed to count how
many copies of $M$ are present in a sample.
But counting requires specifying \emph{where} the count is taken.
In a spatially extended system, the same molecule may be abundant
locally and rare globally.
Assembly Theory provides no spatial or locational framework in which
to ground the count, making copy number undefined in any setting
where location matters.

\paragraph{Gap 2: The identity criterion for copies is undefined.}
Before counting copies of $M$, one must say what it means for two
objects to be copies of the same thing.
Assembly Theory implicitly appeals to structural identity (same
molecular graph), but this criterion is both too strict (ignoring
conformational flexibility) and too informal (not specifying the
relevant equivalence relation).

The Rho calculus fills both gaps simultaneously.

\section{Locations as Channels, Namespaces as Regions}

Among all GSLTs, the Rho calculus is distinguished by providing
an abstract but mathematically precise notion of location.

\begin{definition}[Channel as Location]
In the Rho calculus, a \emph{channel} is a name of the form
$@P$ for a process $P$.
Channels serve as \emph{locations}: the process $P$ is said to
\emph{reside at} the channel $@P$ when it is stored there via
$@P!(P)$ or accessed via $\mathtt{for}(y \leftarrow @P)\,Q$.
Since channels are themselves processes (programs), location is
a computational notion: \emph{where} something is stored is itself
a piece of running code.
\end{definition}

\begin{definition}[Namespace]
A \emph{namespace} is a finite set of channels
$\mathcal{N} = \{@P_1, @P_2, \ldots, @P_k\}$.
A namespace describes a \emph{region} of the computational space:
the collection of locations at which a population is hosted.
\end{definition}

\begin{definition}[Copy Number in a Namespace]
\label{def:copy-number}
Let $P \in \terms(\Rho)$ be a process and let $\mathcal{N}$ be a
namespace.
The \emph{copy number of $P$ in $\mathcal{N}$} is
\[
  \mathrm{CN}(P, \mathcal{N}) \;=\;
  \bigl|\bigl\{ c \in \mathcal{N} \;\mid\; {*}c \bisim P \bigr\}\bigr|
\]
the number of channels in $\mathcal{N}$ whose dereferenced content
is bisimilar to $P$.
\end{definition}

\begin{remark}[Bisimulation as the Identity Criterion]
Definition~\ref{def:copy-number} answers Gap~2 directly: two
processes are copies of each other if and only if they are
\emph{bisimilar} --- experimentally indistinguishable by any
program context.
This is the correct identity criterion for the purposes of the
theory: what matters is not structural identity but behavioral
identity under interaction.
Two molecules are copies in the relevant sense if and only if no
experiment can tell them apart.
\end{remark}

\begin{remark}[The Biosignature Criterion, Made Precise]
With Proposition~\ref{prop:assembly-index} and
Definition~\ref{def:copy-number}, the biosignature criterion of
Assembly Theory becomes the mathematically well-posed statement:
\[
  \mathrm{AI}(P) \gg 1 \quad \text{and} \quad
  \mathrm{CN}(P,\,\mathcal{N}) \gg 1
\]
for some accessible namespace $\mathcal{N}$.
High causal depth (a long construction history) concomitant with
high copy number (abundant presence in a region) is indicative
of a replication process.
\end{remark}

% ===================================================================
